\documentclass[11pt,a4paper]{article}
\usepackage{isabelle,isabellesym}
\usepackage[T1]{fontenc}
\usepackage{amsmath,amssymb}
\usepackage{railsetup}
\usepackage[margin=2cm]{geometry}
\usepackage{url}
\usepackage{hyperref}
\hypersetup{hidelinks,
  pdftitle={Computational Path Presentations: Normal Forms and Changes of Presentation},
  pdfauthor={Arthur Freitas Ramos, David Barros Hulak, Ruy J. G. B. de Queiroz}}
\urlstyle{rm}
\isabellestyle{it}
\title{Computational Path Presentations\\Normal Forms and Changes of Presentation}
\author{Arthur Freitas Ramos \and David Barros Hulak \and Ruy J. G. B. de Queiroz}
\date{13 November 2026}
\begin{document}
\maketitle
\begin{abstract}
Equality of endpoint-typed binary witnesses generated by groupoid equations
is characterized by equality of their reduced signed words. With typed named
equations, flattening is sound and complete for the corresponding word
presentation. Reversible fixed-vertex word interpretations preserve and reflect
equality. Adding derivable equations or defined generators is conservative,
and complete invariants transport across reversible presentations. The
two-loop presentation \(ab=ba\) admits a complete executable normal form
indexed by integer pairs. These are algebraic results; no geometric comparison
or convergence theorem for a full dependent equality rewrite system is claimed.
\end{abstract}
\begingroup
\small
\tableofcontents
\endgroup
\input{session}
\bibliographystyle{abbrv}
\bibliography{root}
\end{document}
